\documentclass[11pt]{article}

\usepackage[utf8]{inputenc}
\usepackage[T1]{fontenc}
\usepackage{lmodern}
\usepackage{geometry}
\usepackage{amsmath,amssymb,amsfonts,amsthm,mathtools}
\usepackage{bm}
\usepackage{enumitem}
\usepackage{microtype}
\usepackage{hyperref}
\usepackage{titlesec}
\usepackage{booktabs}
\usepackage{array}
\usepackage{setspace}
\usepackage{mathrsfs}
\usepackage{physics}

\geometry{margin=1in}
\hypersetup{colorlinks=true, linkcolor=black, urlcolor=blue, citecolor=black}
\setlist[enumerate]{topsep=0.4em,itemsep=0.35em}
\setlist[itemize]{topsep=0.3em,itemsep=0.25em}
\titleformat{\section}{\large\bfseries}{\thesection.}{0.5em}{}
\titleformat{\subsection}{\normalsize\bfseries}{\thesubsection.}{0.5em}{}
\setstretch{1.06}

\newcommand{\Aether}{\AE{}ther}
\newcommand{\Aflow}{\AE{}ther-flow}
\newcommand{\TheoryName}{\Aflow{} Interpretation of Relativity}
\newcommand{\TheoryProgram}{\Aether{} / \Aflow{} framework}
\newcommand{\CompletedMetric}{g^{\sharp}}

\newcommand{\PrePreProtoSectionPullbackCore}{\mathfrak S^{\mathrm{sub,proto,pppresec,pb}}}
\newcommand{\PullbackOperatorCore}{\mathfrak S^{\mathrm{sub,proto,pppresec,pb,op}}}
\newcommand{\PullbackBodyOperatorCore}{\mathfrak S^{\mathrm{sub,proto,pppresec,pb,bodyop}}}
\newcommand{\PullbackBodyResponseCore}{\mathfrak S^{\mathrm{sub,proto,pppresec,pb,bodyresp}}}

\newtheorem{definition}{Definition}
\newtheorem{proposition}{Proposition}
\newtheorem{remark}{Remark}
\newtheorem{corollary}{Corollary}
\newtheorem{theorem}{Theorem}

\title{\textbf{The \TheoryName{}}\\[0.4em]
\large Primitive-Reservoir Observer-Reduction\\
\large Section-Centered Hidden-Response\\
\large Pullback Body-Response Core\\
\large Next Benchmark Criterion}
\author{}
\date{}

\begin{document}

\maketitle

\begin{abstract}
The current active-primary theorem already removed the full observer-relevant pullback body-operator core \(\PullbackBodyOperatorCore\) as the primitive stopping point of the active section-centered hidden-response line: benchmark-facingly, the separate ambient operator-family presentation is now recovered from the smaller observer-relevant pullback body-response core \(\PullbackBodyResponseCore\). Another same-output deeper-origin relay that merely preserves that same body-response core and therefore merely reproduces the same body-operator core, the same pullback-operator core, the same pullback-quadratic core, the same hidden-response package, the same residual-normal-form datum, and the same benchmark-faithful pre-proto-section dynamics no longer counts as the honest default next benchmark-facing manuscript.

This note records the routing consequence of that gain. The next honest benchmark-facing move must either derive or justify the current observer-relevant pullback body-response core \(\PullbackBodyResponseCore\) from still more primitive explicit substrate structure in a way that changes benchmark-facing status, or prove a genuinely smaller theorem-level collapse, sharper necessity result, or sharper uniqueness result strictly inside that smaller core. The benchmark boundary remains fixed throughout. The one operative metric is still exactly \(\CompletedMetric\), the ordinary matter route is unchanged, there is no second operative metric, the low-energy matter law is not enlarged, and no stronger promotion language is licensed. The positive result is therefore routing discipline at the current pullback-body-response-core frontier.
\end{abstract}

\section{Introduction}

The active derivational program of \TheoryName{} is now one level below the pullback-body-operator-core frontier. The preceding theorem already showed that, once the fixed proto-section benchmark base is held fixed, the current pullback-body-operator-core presentation contains no independent benchmark-facing freedom beyond the smaller pullback body-response core \(\PullbackBodyResponseCore\). That theorem therefore spent the full body-operator-core presentation as the live stopping point.

The present note records the routing consequence of that gain. It does not reopen the larger pullback-body-operator-core presentation, the larger pullback-operator-core presentation, the larger pullback-quadratic-core presentation, the larger section-centered hidden-response package, or the already-spent residual-normal-form and pre-proto-section presentations as if they were again independently live. It records only which kind of next manuscript would now count as an honest benchmark-facing gain once the current burden has moved to the smaller body-response core itself.

\paragraph{Framework and claim boundary.}
\TheoryProgram{} denotes the broader \Aether{} / \Aflow{} framework built on the claims that the \Aether{} is the underlying four-dimensional substrate of reality, the \Aflow{} is its intrinsic ordered motion, observed three-dimensional space is the local experiential slice of that deeper substrate, S-time is the experienced order of change arising from matter, light, and the \Aflow{}, and observed expansion is the three-dimensional appearance of deeper four-dimensional ordered motion. Within that framework, \TheoryName{} denotes the exact-closure theory in which the effective gravitational dynamics are adopted to be exactly Einsteinian with universal matter coupling. In the active sequence, \emph{adoption} means use of that established relativistic dynamics without claiming substrate derivation, while \emph{derivation} is reserved for a first-principles recovery from explicit substrate variables.

The present note stays strictly inside that boundary. It records only which kind of next manuscript would now count as an honest benchmark-facing gain at the current pullback-body-response-core frontier.

\section{Fixed Inputs}

\begin{definition}[Pullback-body-response-core frontier inputs]
The criterion recorded here uses the following fixed inputs.
\begin{enumerate}
    \item The benchmark exact-closure package, which fixes the public one-metric GR target of the repository.
    \item The benchmark gatekeeping layer, including the recorded safeguard that blocks same-output deeper-origin relays from counting as new front-line progress once the live benchmark-relevant datum is unchanged.
    \item The earlier theorem proving that the full section-centered hidden-response package carries no independent benchmark-facing freedom beyond the smaller observer-relevant pullback-quadratic core \(\PrePreProtoSectionPullbackCore\).
    \item The later theorem proving that the current pullback-quadratic-core presentation \(\PrePreProtoSectionPullbackCore\) itself carries no independent benchmark-facing freedom beyond the smaller observer-relevant pullback-operator core \(\PullbackOperatorCore\).
    \item The later theorem proving that the current pullback-operator-core presentation \(\PullbackOperatorCore\) itself carries no independent benchmark-facing freedom beyond the smaller observer-relevant pullback body-operator core \(\PullbackBodyOperatorCore\).
    \item The current active-primary theorem proving that the current pullback body-operator-core presentation \(\PullbackBodyOperatorCore\) itself carries no independent benchmark-facing freedom beyond the smaller observer-relevant pullback body-response core \(\PullbackBodyResponseCore\).
\end{enumerate}
\end{definition}

\begin{remark}
The present note does not replace those manuscripts. It records the routing consequence of reading them together under the active safeguard against recursive depth drift.
\end{remark}

\section{Why the Live Burden Now Sits at Pullback Body-Response Scope}

\begin{proposition}[The live burden is renewed at pullback-body-response-core scope]
At the current stage of the primitive-reservoir observer-reduction line, the pullback-body-response-core collapse theorem renews the live benchmark-facing burden at the level of the observer-relevant pullback body-response core \(\PullbackBodyResponseCore\).
\end{proposition}

\begin{proof}
That theorem changes benchmark-facing status because it removes the full pullback-body-operator-core presentation \(\PullbackBodyOperatorCore\) as the primitive stopping point and proves that the current benchmark-facing content of that presentation is already carried by the smaller pullback body-response core \(\PullbackBodyResponseCore\). Once the full ambient operator-family presentation is shown to carry no independent benchmark-facing freedom beyond that smaller core, the live burden no longer sits on the body-operator-core presentation itself. It sits on the sharper body-response core that now determines the same body-operator core, the same pullback-operator core, the same pullback-quadratic core, and the same downstream benchmark-facing consequences at current scope.
\end{proof}

\begin{definition}[Benchmark-neutral deeper relay at current pullback-body-response-core scope]
A \emph{benchmark-neutral deeper relay at current pullback-body-response-core scope} is a manuscript that preserves the same benchmark one-metric output, the same ordinary matter route, the same current pullback body-response core \(\PullbackBodyResponseCore\), the same current pullback body-operator core \(\PullbackBodyOperatorCore\), the same current pullback-operator core \(\PullbackOperatorCore\), the same current pullback-quadratic core \(\PrePreProtoSectionPullbackCore\), the same current hidden-response package consequences, the same current residual-normal-form datum, the same current benchmark-faithful pre-proto-section package, and the same downstream benchmark-facing consequences while merely relocating the unexplained substrate layer deeper, repackaging the same smaller core, or re-expanding the presentation back to the larger pullback-body-operator-core, pullback-operator-core, pullback-quadratic-core, hidden-response, residual-normal-form, or pre-proto-section presentations without a new compression, necessity, uniqueness, or comparably sharp routing gain.
\end{definition}

\section{The Current Pullback Body-Response Criterion}

\begin{definition}[Admissible next benchmark-facing manuscript at current pullback-body-response-core scope]
At the current pullback-body-response-core frontier, a manuscript counts as an admissible next benchmark-facing move only if it does at least one of the following:
\begin{enumerate}
    \item derives or justifies the current observer-relevant pullback body-response core \(\PullbackBodyResponseCore\) from still more primitive explicit substrate structure in a way that does more than merely relocate the unexplained layer deeper;
    \item proves a genuinely smaller theorem-level collapse, sharper necessity result, or sharper uniqueness result strictly inside \(\PullbackBodyResponseCore\);
    \item does either of the previous two while preserving the same benchmark one-metric package, the same ordinary matter route, and the same claim boundary.
\end{enumerate}
\end{definition}

\begin{theorem}[Next honest benchmark-facing task at current pullback-body-response-core scope]
The next honest benchmark-facing manuscript below the current pullback-body-response-core frontier must either derive or justify the current observer-relevant pullback body-response core \(\PullbackBodyResponseCore\) from still more primitive explicit substrate structure in a benchmark-relevant way, or else prove a genuinely smaller collapse, sharper necessity result, or sharper uniqueness result strictly inside that smaller core. A manuscript that only repeats the same current body-response datum through another deeper relay, re-expands the discussion back to the larger pullback-body-operator-core, pullback-operator-core, pullback-quadratic-core, hidden-response, residual-normal-form, or pre-proto-section presentations without such a new gain, or invokes a quantum branch without doing the same benchmark-facing work does not satisfy the criterion.
\end{theorem}

\begin{proof}
By Proposition 1, the live burden already lies at the smaller pullback body-response core \(\PullbackBodyResponseCore\) rather than at the larger pullback-body-operator-core presentation \(\PullbackBodyOperatorCore\). Therefore a subsequent manuscript must improve on that smaller core rather than merely accompany it. A deeper-origin manuscript can do so only if it changes benchmark-facing status by more than depth alone. If it simply reproduces the same current body-response core and its same downstream outputs, the routing safeguard classifies it as benchmark neutral. The remaining honest alternatives are therefore exactly those stated in the definition: either a more primitive derivation or justification of the current body-response core itself, or a sharper internal collapse, necessity result, or uniqueness result strictly inside it.
\end{proof}

\section{What the Next Move Should Not Be}

\begin{proposition}[Screened next moves at the current pullback-body-response-core frontier]
At the current pullback-body-response-core frontier, none of the following is the default next benchmark-facing move:
\begin{enumerate}
    \item another same-output deeper-origin relay that merely preserves the current pullback body-response core \(\PullbackBodyResponseCore\) and therefore merely reproduces the same current pullback body-operator core \(\PullbackBodyOperatorCore\), the same current pullback-operator core \(\PullbackOperatorCore\), the same current pullback-quadratic core \(\PrePreProtoSectionPullbackCore\), the same current hidden-response package, the same current residual-normal-form datum, the same current benchmark-faithful pre-proto-section package, and the same downstream benchmark-facing consequences;
    \item another re-expansion back to the larger pullback-body-operator-core presentation, to the larger pullback-operator-core presentation, to the larger pullback-quadratic-core presentation, to the larger hidden-response package, to the already-spent residual-normal-form presentation, or to the larger pre-proto-section presentation as if those larger presentations were again independently live burdens;
    \item a quantum branch that does not derive or justify the same body-response core, collapse it further, or close a benchmark derivation gate in a genuinely new way;
    \item a manuscript that introduces a second operative metric, enlarges the ordinary matter law, or uses stronger promotion language.
\end{enumerate}
\end{proposition}

\begin{proof}
Item (1) fails the preceding criterion because it leaves the renewed live burden unchanged. Item (2) likewise fails because it re-expands to larger already-spent presentations rather than removing remaining freedom in \(\PullbackBodyResponseCore\) or proving a smaller internal datum there. Item (3) does not directly address the live burden at current scope unless it performs the same benchmark-facing work named above. Item (4) violates the fixed claim boundary of the benchmark package and therefore cannot count as a disciplined next step on the active line. The benchmark package remains the exact one-metric GR package already fixed by adoption, so any admissible next manuscript must preserve that boundary.
\end{proof}

\begin{remark}
This screening statement is not a denial that those deeper or alternative manuscripts may matter strategically. It is a routing statement: they are not the default way to spend the next benchmark-facing move from the current pullback-body-response-core frontier unless they do the same benchmark-facing work named above.
\end{remark}

\section{Conclusion}

The positive result of this note is routing discipline at the current pullback-body-response-core frontier. The pullback-body-response-core collapse theorem remains the live benchmark-facing gain because it already removes the full pullback-body-operator-core presentation as the primitive stopping point by proving that the current benchmark-facing content of that presentation is carried by the smaller observer-relevant pullback body-response core. The earlier pullback-body-operator-core theorem remains preserved main-line history, but under the recursive-depth safeguard it is no longer itself the default current frontier.

The scope remains narrow and honest. The benchmark package is unchanged: the one operative metric is still exactly \(\CompletedMetric\), the ordinary matter route is unchanged, there is no second operative metric, and the low-energy matter law is not enlarged. Nothing here upgrades adoption to derivation or licenses stronger promotion language.

The next open burden is therefore clear. The next benchmark-facing manuscript should derive or justify the current pullback body-response core from still more primitive explicit substrate structure in a way that changes benchmark-facing status, unless the sharper honest gain available instead is a genuinely smaller theorem-level collapse, sharper necessity result, or sharper uniqueness result strictly inside that smaller core. Another same-output deeper relay that merely preserves that body-response core, a re-expansion back to the larger pullback-body-operator-core, pullback-operator-core, pullback-quadratic-core, hidden-response, residual-normal-form, or pre-proto-section presentations, or a quantum branch that does not do that same benchmark-facing work is not the clean next manuscript target at current scope.

\end{document}
